\documentclass[10pt,a4paper]{amsart}
\usepackage{amsmath,amssymb,amsthm,graphicx}
\usepackage{verbatim}
\newenvironment{reference}{\comment}{\endcomment}
%\newenvironment{reference}{}{}
\newenvironment{slogan}{\comment}{\endcomment}
\newenvironment{history}{\comment}{\endcomment}

% For commutative diagrams we use Xy-pic
\usepackage[all]{xy}

% We use 2cell for 2-commutative diagrams.
\xyoption{2cell}
% Theorem environments.
%
\theoremstyle{plain}
\newtheorem{theorem}[subsection]{Theorem}
\newtheorem{proposition}[subsection]{Proposition}
\newtheorem{lemma}[subsection]{Lemma}

\theoremstyle{definition}
\newtheorem{definition}[subsection]{Definition}
\newtheorem{example}[subsection]{Example}
\newtheorem{exercise}[subsection]{Exercise}
\newtheorem{situation}[subsection]{Situation}

\theoremstyle{remark}
\newtheorem{remark}[subsection]{Remark}
\newtheorem{remarks}[subsection]{Remarks}

\numberwithin{equation}{subsection}

% Macros
%
\def\lim{\mathop{\mathrm{lim}}\nolimits}
\def\colim{\mathop{\mathrm{colim}}\nolimits}
\def\Spec{\mathop{\mathrm{Spec}}}
\def\Hom{\mathop{\mathrm{Hom}}\nolimits}
\def\Ext{\mathop{\mathrm{Ext}}\nolimits}
\def\SheafHom{\mathop{\mathcal{H}\!\mathit{om}}\nolimits}
\def\SheafExt{\mathop{\mathcal{E}\!\mathit{xt}}\nolimits}
\def\Sch{\mathit{Sch}}
\def\Mor{\mathop{\mathrm{Mor}}\nolimits}
\def\Ob{\mathop{\mathrm{Ob}}\nolimits}
\def\Sh{\mathop{\mathit{Sh}}\nolimits}
\def\NL{\mathop{N\!L}\nolimits}
\def\CH{\mathop{\mathrm{CH}}\nolimits}
\def\proetale{{pro\text{-}\acute{e}tale}}
\def\etale{{\acute{e}tale}}
\def\QCoh{\mathit{QCoh}}
\def\Ker{\mathop{\mathrm{Ker}}}
\def\Im{\mathop{\mathrm{Im}}}
\def\Coker{\mathop{\mathrm{Coker}}}
\def\Coim{\mathop{\mathrm{Coim}}}

% Boxtimes
%
\DeclareMathSymbol{\boxtimes}{\mathbin}{AMSa}{"02}

%
% Macros for moduli stacks/spaces
%
\def\QCohstack{\mathcal{QC}\!\mathit{oh}}
\def\Cohstack{\mathcal{C}\!\mathit{oh}}
\def\Spacesstack{\mathcal{S}\!\mathit{paces}}
\def\Quotfunctor{\mathrm{Quot}}
\def\Hilbfunctor{\mathrm{Hilb}}
\def\Curvesstack{\mathcal{C}\!\mathit{urves}}
\def\Polarizedstack{\mathcal{P}\!\mathit{olarized}}
\def\Complexesstack{\mathcal{C}\!\mathit{omplexes}}
% \Pic is the operator that assigns to X its picard group, usage \Pic(X)
% \Picardstack_{X/B} denotes the Picard stack of X over B
% \Picardfunctor_{X/B} denotes the Picard functor of X over B
\def\Pic{\mathop{\mathrm{Pic}}\nolimits}
\def\Picardstack{\mathcal{P}\!\mathit{ic}}
\def\Picardfunctor{\mathrm{Pic}}
\def\Deformationcategory{\mathcal{D}\!\mathit{ef}}
\title{A two-sided subspace of a tensor product is generated by its intersection with the constants}
\author{The Stacks Project Authors}
\date{stacks-project:brauer.tex:local-snapshot-20261012; GNU FDL 1.2 or later}
\begin{document}
\maketitle
\noindent\textit{Editorial note.} All delivered mathematical text is the author's own, copied from the pinned Stacks Project chapter file sites-cohomology.tex; no statement or proof step was rewritten, reordered or removed.\par\medskip
\noindent\textit{Editorial note.} The proof is the author's complete argument, both the reduction to the quotient configuration and the minimality argument being written out with the elements displayed; no step is left to the reader.\par\medskip
\noindent\textit{Editorial note.} Every labelled reference occurring inside the counted spans is carried in assembly.external\_references as a non-load-bearing navigation pointer (same-chapter ones at their definition site, cross-chapter ones at the invocation inside the counted span), so no load-bearing mathematics is replaced by a stub. No bibliography citation occurs inside the counted spans.\par\medskip
\noindent\textit{Editorial note.} The same-name check for this card scanned every staged Stacks chapter of this repository -- all sources/stacks-*/pinned-original-source.tex files, the card's own chapter excluded -- and the labels found are recorded in dedup.same\_name\_cross\_chapter on the card.\par\medskip
\noindent\textit{Editorial note.} The preamble is the chapter's own retained minimal preamble (sources/stacks-sites-cohomology/retained-preamble.tex), a verbatim excerpt of the native preamble: the Xy-pic setup and the native theorem declarations with the mathematical macros.\par\medskip
\noindent\textit{Allowed standard prerequisites (editorial).}\begin{itemize}
\item Central algebras over a field. Same chapter, the definition of a central algebra and the identification of its center with the ground field. Supplies that the center of the skew field is exactly the ground field, which is what forces the coefficients of a minimal expression to lie in the ground field.
\item Tensor products of vector spaces. Same chapter, the elementary properties of the tensor product of a vector space with an algebra. Supply the identification used to pass to the quotient by the subspace generated by the elements whose constant tensor lies in the given subspace.
\end{itemize}\medskip
\begin{lemma}
\label{lemma-generate-two-sided-sub}
Let $V$ be a $k$ vector space. Let $K$ be a central $k$-algebra
which is a skew field. Let $W \subset V \otimes_k K$ be a two-sided
$K$-sub vector space. Then $W$ is generated as a left $K$-vector
space by $W \cap (V \otimes 1)$.
\end{lemma}

\begin{proof}
Let $V' \subset V$ be the $k$-sub vector space generated by $v \in V$
such that $v \otimes 1 \in W$. Then $V' \otimes_k K \subset W$ and
we have
$$
W/(V' \otimes_k K)  \subset  (V/V') \otimes_k K.
$$
If $\overline{v} \in V/V'$ is a nonzero vector such that
$\overline{v} \otimes 1$ is contained in $W/(V' \otimes_k K)$,
then we see that $v \otimes 1 \in W$ where $v \in V$ lifts $\overline{v}$.
This contradicts our construction of $V'$. Hence we may replace
$V$ by $V/V'$ and $W$ by $W/(V' \otimes_k K)$ and it suffices to prove
that $W \cap (V \otimes 1)$ is nonzero if $W$ is nonzero.

\medskip\noindent
To see this let $w \in W$ be a nonzero element which can be written
as $w = \sum_{i = 1, \ldots, n} v_i \otimes k_i$ with $n$ minimal.
We may right multiply with $k_1^{-1}$ and assume that $k_1 = 1$.
If $n = 1$, then we win because $v_1 \otimes 1 \in W$.
If $n > 1$, then we see that for any $c \in K$
$$
c w - w c = \sum\nolimits_{i = 2, \ldots, n} v_i \otimes (c k_i - k_i c) \in W
$$
and hence $c k_i - k_i c = 0$ by minimality of $n$.
This implies that $k_i$ is in the center of $K$ which is $k$ by
assumption. Hence $w = (v_1 + \sum k_i v_i) \otimes 1$ contradicting
the minimality of $n$.
\end{proof}

\end{document}
